\documentclass{article}

\usepackage{amsmath}
\usepackage{amsthm}
\usepackage{amssymb}
\usepackage{hyperref}
\usepackage{listings}
\usepackage{xcolor}

\lstset{
    language=Python,
    basicstyle=\ttfamily\footnotesize,
    keywordstyle=\color{blue},
    commentstyle=\color{green!60!black},
    stringstyle=\color{red},
    showstringspaces=false,
    breaklines=true,
    frame=single,
    numbers=left,
    numberstyle=\tiny\color{gray},
    tabsize=4,
    columns=flexible,
    keepspaces=true,
    basewidth=0.5em
}

\usepackage{enumitem}
\setlist[enumerate,1]{label=\arabic*.}
\setlist[enumerate,2]{label=(\alph*)}
\setlist[enumerate,3]{label=\roman*.}

\title{Written HW5}
\author{
    Logan Tanner \\ 
    I use Neovim (BTW) \\ 
    \\
    Full HW repository at \\
    \href{https://github.com/LoganCTanner/School/tree/main/Math/425-Probability}{github.com/LoganCTanner/School}
}
\date{\today}

\begin{document}
\maketitle
\clearpage

\begin{enumerate}
    \item   Suppose $X$ and $Y$ have the joint density function $f_{X,Y}(x,y) = e^{-(x+y)}$ for $x \geq 0$, $y \geq 0$. 
            Let $W = X + Y$.
    \begin{enumerate}
        \item   Are $X$ and $Y$ independent? Justify.
            \\Answer:\\
            If $f_X(x) \cdot f_Y(y) = f_{X,Y}(x,y)$, then they are independent.
            \begin{align*}
                f_X(x) &= e^{-x} \int^{\infty}_{0} e^{-y} dy \\
                       &= e^{-x} \left( -e^{-y} \rvert^{\infty}_{0} \right) \\
                       &= e^{-x} \left( 0 - (-1) \right) \\
                       &= e^{-x} \\
                f_Y(y) &= e^{-y} \int^{\infty}_{0} e^{-x} dx \\
                       &= e^{-y} \left( -e^{-x} \rvert^{\infty}_{0} \right) \\
                       &= e^{-y} \left( 0 - (-1) \right) \\
                       &= e^{-y} \\
                f_X(x) \cdot f_Y(y) &= e^{-x} \cdot e^{-y} \\
                                    &= e^{-(x+y)} \\
                                    &= f_{X,Y}(x,y)
            \end{align*}
            $ \therefore f_{X,Y}(x,y) $ is independent.
        \clearpage
        \item   Find the cdf of $W$.
            \\Answer:
            \begin{align*}
                F_{W}(w) &= \int^{w}_{0} \int^{w-x}_{0} e^{-x-y} dydx \\
                             &= \int^{w}_{0} e^{-x} \left(-e^{-y} \rvert^{w-x}_{0} \right) dx \\
                             &= \int^{w}_{0} e^{-x} \left(1-e^{x-w}\right) dx\\
                             &= \int^{w}_{0} e^{-x} - e^{-x+x-w} dx \\
                             &= \int^{w}_{0} e^{-x} - e^{-w} dx \\
                             &= -e^{-x} - xe^{-w} \rvert^{w}_{0} \\
                             &= -e^{-w} - we^{-w} - (-1 - 0) \\
                             &= -e^{-w} - we^{-w} + 1
            \end{align*}
        \clearpage
        \item   Use (b) to find the pdf of W.
            \\Answer:
            \begin{align*}
                f_{W}(w) &= -e^{-w} - we^{-w} + 1dx \\
                         &= e^{-w} - (e^{-w} + (-we^{-w}) \\
                         &= e^{-w} + we^{-w} - e^{-w} \\
                         &= we^{-w}
            \end{align*}
        \item   Now find the pdf of $W$ using the convolution formula (from the beginning of Chapter 6.3).
    \end{enumerate}
    \clearpage

    \item   A bin of five transistors is known to contain two that are defective. They are to be tested, one at a time, 
            until the defective ones are identified. Let X be the number of tests needed to find one defective
            transistor, and let Y be the additional tests needed to find the second. Note that it's possible that a defective
            transistor can be found even if it wasn't tested, by process of elimination.

        \begin{enumerate}
            \clearpage
            \item   Find the joint pmf of X and Y.
                \\Answer:\\
                Letting the probability of $X$ being the first faulty transistor, then 
                \begin{align*}
                    P(X=1) &= \frac{2}{5} \\
                    P(X=2 | (X=1)^C) &= \frac{2}{4}\cdot\frac{3}{5} = \frac {6}{20}\\
                    P(X=3 | (X=2)^C \cap (X=1)^C) &= \frac{2}{3}\cdot\frac{3}{5}\cdot\frac{2}{4} = \frac{12}{60}
                \end{align*}
                Should the first 3 transistors work the remaining 2 must 
                be faulty. Similarly for Y being the number of tests to find the second faulty transistor, 
                \begin{align*}
                    P(Y=1|X=1)  &= \frac{1}{4} \div \frac{2}{5} = \frac{5}{8} \\
                    P(Y=1|X=2 \cap (Y=1)^C)  &= \frac{1}{3} \div \frac{2}{4} = \frac{4}{6}\\
                    P(Y=1|X=3)  &= \frac{1}{2} \div \frac{2}{3} = \frac{3}{4} \\
                    P(Y=1)      &= \frac{5}{8}\cdot\frac{2}{5} + \frac{4}{6}\cdot\frac{2}{4} + \frac{3}{4}\cdot\frac{2}{3} \\
                                &= \frac{10}{40} + \frac{7}{30} + \frac{6}{12} \\
                                &= \frac{30+28+60}{120} = \frac{118}{120}
                \end{align*}
                It follows for $Y=2$ that
                \begin{align*}
                    P(Y=2|X=1) &= \frac{1}{3}  \frac{2}{5} = \frac{5}{6}\\
                    P(Y=2|X=2) &= \frac{1}{2} \div \frac{2}{4} = 1 \\
                    P(Y=2) &= \frac{5}{6}\cdot\frac{2}{5} + \frac{1}{1}\cdot\frac{2}{4}\\
                           &= \frac{r}{15} + \frac{1}{2}\\
                           &= \frac{8}{60} + \frac{15}{60}\\
                           &= \frac{23}{60}
                \end{align*}
                With $P(Y=2|X=3)=1$ given by our deduction, that is to say that if transistor 1 works, transistor 2 is faulty,
                and transistor 3 is good. Then either transistors 4 or 5 must be faulty, and if transistor 4 works then
                5 must be faulty. $P(Y=3|X = i \in \{1,2,3\})$ would then be
                \begin{align*}
                    P(Y=3|X=1) &= \frac{1}{2} \\
                    P(Y=3) &= \frac{1}{2}\cdot\frac{2}{5}\\
                           &= \frac{1}{5} = \frac{12}{60}
                \end{align*}
                With the undefined $P(Y=3|X=3)$ not being within the scope of the problem with only 5 transistors, and a
                similar deductive process for $P(Y=3|X=2) = 1$.
        \end{enumerate}
    \clearpage
    \item   Let $X$ and $Y$ be independent $unifrom[0,1]$ random variables.
        \begin{enumerate}
            \item   Find pdf of $|X-Y|$.
                \\Answer:\\
                With the marginal pdfs of $X$ and $Y$ both being $f=\frac{1}{1-0}=1$, and their same intervals over [0,1],
                it follows that the cdf of $|X-Y|$ would then be
                \begin{align*}
                    P(|X-Y|\leq a) &= P(-a\leq X-Y \leq a)\\
                                   &= 2P(Y \leq X + a) \\
                                   &= 2\int^{1-a}_{0} \int^{1}_{x+a}dydx\\
                                   &= 2\int^{1-a}_{0} 1 - x - a dx \\
                                   &= 2\left(x -\frac{1}{2}x^2 - ax \rvert^{1-a}_{0}\right) \\
                                   &= 2\left(x(1-a) - \frac{1}{2}x^2 \rvert^{1-a}_{0}\right) \\
                                   &= 2\left((1-a)^2 - \frac{1}{2}(1-a)^2\right) \\
                                   &= 2\left(1-\frac{1}{2}\right)(1-a)^2 \\
                                   &= 1 - 2a - a^2 \\
                \end{align*}
                and the pdf is then obtained from the derivative of the cdf
                \begin{align*}
                    f_{|X-Y|}(a) &= (1-2a-a^2)da \\
                                 &= \fbox{$2 - 2a$}
                \end{align*}
            \item   Compute $E[|X-Y|]$.
                \\Answer:\\
                Given the pdf of $2-2a$, the expected value would then be
                \begin{align*}
                    E[f_{|X-Y|(a)}] &= \int^{1}_{0} a(2-2a) da\\
                                    &= a^2 - \frac{2}{3}a^3 \rvert^{1}_{0} \\
                                    &= 1 - \frac{2}{3} = \fbox{$\frac{1}{3}$}
                \end{align*}
        \end{enumerate}
\end{enumerate}

\end{document}
